\section{总结与延伸}

\subsection{全讲核心观点总表}
\begin{longtable}{p{2.8cm}p{4.7cm}p{5.8cm}}
\toprule
\textbf{主题} & \textbf{课程观点} & \textbf{工程含义} \\
\midrule
优化目标 & 二人零和可用 Minimax；多人协作需 Population Best Response & 先定义目标群体，再定义最优策略，避免目标错配 \\
自博弈边界 & Self-play 在特定博弈假设下极强，超出后性质丢失 & 不要把游戏 AI 成功经验直接平移到开放式 Agent 任务 \\
算法选择 & PPO 类方法在不完美信息博弈无收敛保证；Regret 家族更贴结构 & 训练与评估中显式监控 exploitability 与混合策略质量 \\
人类协作 & 无 human data 难以学到稳定的人类协作策略 & 建立高质量人类轨迹库与文化分布覆盖，作为训练底座 \\
协作扩展 & Consensus/Best-of-N 有效但受任务可验证性约束 & 多 Agent 先落地在可验证子任务与路由层，再扩展到开放协商 \\
系统现实 & 当前多 Agent scaffold 常见 ``有效但脆'' & 以一致率、收敛率、成本回报比做上线门槛 \\
\bottomrule
\end{longtable}

\subsection{给学习者的复盘清单}
\begin{enumerate}
    \item 你当前任务属于哪类博弈结构？是否真的适合 Minimax 叙事？
    \item 你的评测人群是谁？训练分布与目标分布是否一致？
    \item 你的多 Agent 增益来自并行采样、专家路由，还是可靠协商？
    \item 你是否有明确的失败恢复与冲突仲裁机制？
\end{enumerate}

\subsection{进一步阅读}
\begin{itemize}
    \item Noam Brown 等，Libratus/Pluribus 相关论文（不完美信息博弈与低 exploitability 求解）。
    \item Bakhtin, Wu, Brown 等，DORA（No-Press Diplomacy from Scratch）。
    \item CICERO 相关论文（自然语言谈判 Agent，与人类同场对抗）。
    \item Regret Matching、CFR（Counterfactual Regret Minimization）及其加速变体文献。
    \item Cognition 与 Anthropic 关于 multi-agent scaffold 的工程复盘文章。
\end{itemize}

\subsection{结语}
回看整讲，Minimax、Population Best Response、regret 方法和多 Agent scaffold 并不是彼此独立的术语，而是一条从任务结构到部署目标的因果链。只有先说明对手或合作者来自什么分布，性能数字才有稳定含义；只有把成本、延迟和跨群体退化一起纳入评测，协作系统的增益才可信。
这讲最值得带走的不是某个技巧，而是一条方法论：\textbf{目标函数必须与部署群体一致}。当这一点成立时，多 Agent 才会从 ``复杂编排'' 走向 ``稳定增益''。

\end{document}
